\chapter{Discussion}
\label{chap:discussion}

Chapters 5 and 6 leave a clear pattern. House-purchase lending growth shows the most durable positive response to an ECB easing surprise, the two lending spreads move with it, the compensation-linked variables respond less consistently, and the real DAX falls. \Cref{tab:persistence-synthesis,tab:hypothesis-assessment} summarize that pattern. This chapter asks what it means. Each section takes one block of the evidence, reads it as economics, and sets it against what earlier studies found. The reading stays reduced-form and comparative throughout. The chapter closes with the limits of the design and the research agenda that follows from them.

\section{Housing-Credit Transmission as the Core Result}

The main result is that ECB policy-news surprises leave their clearest and most persistent trace in house-purchase lending growth. The response is positive on impact, peaks at six months, and stays positive through the first year. The single-horizon estimate turns negative at two years, so part of the first-year accumulation reverses, and the accumulated response falls from about 0.48 to about 0.36 while staying positive. Its HAC interval excludes zero at two years, the dependence-aware bootstraps do not, and all three agree that the accumulated response excludes zero through the first year. Sign stability across the event and exclusion samples is high, the rolling and recursive windows keep the positive sign in most cases, and the direction of the accumulated path never changes. No other variable combines an early response, a durable accumulation, and this degree of sign coherence.

The contrast with the rest of the credit block is what makes the result specific. Broad NFC and household credit stocks barely move, and their accumulated intervals include zero at every horizon. Easing does not show up as a general expansion of credit. It shows up in the one credit aggregate that finances a collateral-intensive purchase. That is why the housing-credit finding is the center of the thesis rather than a supporting detail inside an asset-price story, and why it is read as a channel-specific statement about transmission and not as a statement about house prices or borrower access, which sit outside the design.

Earlier studies point the same way. \textcite{BernankeGertler1995} report, for the United States, that residential investment is the spending component that responds most sharply to a monetary policy shock, and they read that as the signature of the credit channel. They also caution that credit quantities alone cannot show that a credit channel is at work, because quantities mix supply and demand. \textcite{JordaSchularickTaylor2015} find in a long annual panel of advanced economies that loose monetary conditions feed booms in real estate lending and house prices, and \textcite{JordaSchularickTaylor2016} show that mortgage lending has grown into the largest part of bank lending. The present result agrees with all three and adds a piece they do not measure. Their identification rests on recursive VAR innovations in the first case and on the policy trilemma in the 2015 panel study, and their housing measures are spending or credit aggregates. Here the same channel is the clearest response to a high-frequency ECB surprise, traced month by month over two years, in a euro-area sample that runs from 2005 to 2025 and covers the financial crisis, the low-rate and QE years, the pandemic, and the tightening that followed.

\section{Stocks, Flows, and the Lending Result}

The housing-credit result rests on a stock-based growth rate, and the choice matters for how it should be read. House-purchase lending growth measures the growth of the outstanding book. It captures cumulative financing conditions and sustained accommodation over the medium run. Pure-new mortgage loans measure monthly originations. They capture immediate lending behavior and are exposed to a short sample, to refinancing composition, and to month-specific reporting noise.

The two series behave as that distinction predicts. The growth of the stock gives the strong and persistent response. New loans are weaker and less stable. Their sample is short, their sign stability is close to a coin toss, and their peak interval includes zero under HAC standard errors even where the bootstrap band is tighter. There is no contradiction in this. Sustained credit accommodation appears in the growth of the housing-credit stock before it appears cleanly in monthly new flows. The two series are also linked in time. An early move in the growth of the stock is associated with a later move in new flows, and the coefficient on that link is positive and precisely estimated. The flow series therefore stays a robustness surface rather than a replacement for the canonical proxy.

The distinction also fixes which literature the result should be compared with. The housing-credit-cycle evidence in \textcite{JordaSchularickTaylor2016} is built on outstanding mortgage credit, a stock, so the durable stock-based response is the closer counterpart to that work. The noisier flow series has no clean analogue there and is not asked to carry a comparison.

\section{Financial Intermediation Mechanism}

The intermediation evidence is what makes the housing-credit result economically believable. The chain in \Cref{eq:mechanism-chain} runs from policy news to lending conditions, then to housing credit, and finally to the real-economy block. The lending-condition step is the one that has to hold for the chain to mean anything, and it does. The mortgage and NFC spreads over the deposit facility rate widen after an easing surprise, and their accumulated responses stay positive through two years. The mortgage spread is the most stable response in the whole set, with a perfect sign-stability score and an accumulated interval that excludes zero from one month on under the HAC and the moving-block intervals.

The direction of that move needs one line of explanation, because it is the one place where the thesis appears to conflict with the credit-cost literature. \textcite{GertlerKaradi2015} find that modest surprise movements in the policy rate produce large movements in credit costs, mainly through term premia and credit spreads, and conclude that the standard interest-rate channel misses part of transmission. Their spreads are private borrowing rates over risk-free rates of matching maturity, so a rise means dearer credit. The spreads here are retail lending rates over the policy rate. They widen after an easing surprise while the lending rates themselves show no precise response, so the margin over the policy rate opens. That points to incomplete pass-through rather than a rise in the cost of credit relative to the risk-free curve, and the two findings are complementary. Both say that the policy rate alone does not describe credit conditions and that the intermediary sits between policy and the borrower. The sustained margin matches the intermediary view in \textcite{AdrianShin2010} and \textcite{BorioZhu2012}, in which policy works through bank balance sheets and loan pricing before it works through volumes. Bank-level euro-area evidence adds that pass-through is uneven. \textcite{AltavillaCanovaCiccarelli2020} show that pass-through from policy rates to lending rates differs across banks with their capital ratios and sovereign exposures, and that the non-standard measures compressed lending rates most at the weakest banks.

The credit-quantity step is where a naive story breaks. If easing simply expanded all credit, broad NFC and household stocks would move with it. They do not. The picture is sectorally differentiated transmission, not a generic expansion of credit quantities, and it is the pattern the credit-channel theory in \Cref{chap:theory} predicts for a collateral-intensive sector.

What the timing outputs add is more modest than the chain might suggest. The sequential pathways in \Cref{tab:sequential-transmission} fix the upstream variable at an early horizon and the downstream variable at a later one, so their order is imposed rather than discovered. Within that design the quantity-to-quantity links are positive and precise, while the link from an early widening of the mortgage margin to later house-purchase growth is negative and marginal. A wider margin in the first month goes with somewhat weaker housing-credit growth six months later, which is the sign a price-quantity relation should carry. The peak horizons tell the same story from another angle. House-purchase lending growth peaks at six months and the spreads at twelve, so the spreads do not lead housing credit. The evidence for intermediation is therefore co-movement on the same shock, a sustained margin and a sustained quantity response over the same year, rather than an observed sequence from one to the other. Monthly aggregates could not separate one bank's decision from one household's response even if the sequence were visible, so the thesis describes a financial propagation with a price side and a quantity side and stops there.

\section{Equity Responses and Information Effects}

The DAX is the hardest financial result to read. It is negative after an easing surprise, and the negativity deepens from three months to two years. The twelve- and twenty-four-month intervals exclude zero at both the 90 and 95 percent levels, and the accumulated path is durable with a perfectly consistent cumulative direction. A response this stable has to be interpreted rather than set aside.

Forcing it into an asset-appreciation story would be wrong, and the literature explains why. An ECB announcement carries two things at once, the policy accommodation and the central bank's own reading of the economy. \textcite{JarocinskiKaradi2020} separate the two by the co-movement of interest rates and stock prices inside the announcement window. A pure policy tightening raises rates and lowers equities, while a positive central-bank information shock raises both, and they show that ignoring the information component biases estimates of what policy does. \textcite{NakamuraSteinsson2018} reach the same conclusion for the Federal Reserve from the reaction of expected output growth to policy surprises. When the information component of an easing surprise is adverse, expected earnings fall, risk premia rise, and the equity market can drop even as policy loosens.

That is the situation the event screen was built for, and the screen says the situation is common. It classifies 132 of the 208 events in the identification window as contaminated or ambiguous, so a large share of the raw target surprises may carry an information component. A negative equity response to a surprise signed as easing is then what the literature predicts, not a contradiction of it. The DAX marks the limit of the result rather than its foundation. The liquidity stock fits the same reading, negative but imprecise, while the housing-credit and lending-condition blocks hold the mechanism.

\section{Compensation-Linked Responses}

The compensation block is weaker, not absent, and that distinction is what the second hypothesis turns on. The real wage tracker excluding one-offs moves at short and intermediate horizons, and its accumulated response excludes zero at six and twelve months under the HAC and wild intervals. By two years that exclusion is gone under every method. The series loses precision exactly where the housing-credit result keeps it. Employment expectations tell a similar story from the other side. They respond early and by a lot, then reverse to a negative and individually significant value at two years. An early reaction that unwinds is not durable income transmission.

The robustness proxies qualify this. The German industry wage-bill measure gives a more positive and durable signal, but it is a national and sectoral series, and \Cref{chap:data} sets out why it stays a secondary proxy.

The weak response is also what earlier work would lead one to expect. \textcite{ChristianoEichenbaumEvans1999} summarize the VAR evidence for the United States as follows: after a monetary policy shock real wages move only by a small amount, and the price level adjusts with a substantial delay. The euro-area wage tracker also moves only a little after a high-frequency easing surprise, with a short-lived real dip and no precise long-horizon movement. It is also what the staggered-contract logic behind \Cref{hyp:compensation-linked-attenuation} predicts, since negotiated wages can respond to a surprise only as agreements are renewed \parencite{Taylor1980}. What the thesis adds is the comparison. The same shock produces a durable housing-credit response and a fading compensation response, which turns a familiar observation about wages into a statement about channels.

The compensation result is therefore comparative. The euro-area wage tracker and employment expectations respond less consistently and less persistently than house-purchase lending growth and the lending spreads. The German wage bill does not. Whether income gains were absent across households is a different question, one that needs microdata and a different identification design, and \Cref{sec:future-research} records it.

\section{Post-COVID and Regime Interpretation}

The regime outputs describe historical dependence and do not replace the full-sample estimates. The reason is mechanical. The thesis is about persistence, and persistence needs long horizons. The COVID window has too few observations for any horizon beyond one month, and the tightening regime is short relative to the two-year horizon. A short window cannot identify a two-year response profile, so the regime evidence is kept descriptive by design. This is also the main limit on what the thesis can say about the post-COVID period itself. The months since 2020 are under a third of the sample, so the full-sample estimates describe how ECB policy news transmits across 2005 to 2025, with the post-COVID cycle as one part of that record rather than a separately estimated window.

Read that way, it still adds something. The tightening regime shows positive housing-credit and spread accumulation with negative wage-pressure accumulation, which lines up with the full-sample ordering. The COVID window shows large moves at the impact and one-month horizons across employment expectations, wage pressure, the DAX, and house-purchase lending growth, which may signal state dependence. The pre-QE and QE windows map the same target surprise into very different observed conditions, with strongly positive accumulated house-purchase growth before QE and negative house-purchase growth and NFC spreads under it. The size, and sometimes the sign, of a response can depend on the regime. The full-sample estimates remain the primary result, and the regime evidence bounds how far it should be generalized.

\section{Implications for Research and Practice}

The findings have moderate but real relevance for practice, and the relevance lies in the specific numbers rather than in the general point that channels differ. Three implications follow. First, the horizon at which transmission is assessed changes the verdict. The housing-credit response peaks at six months, its accumulated value peaks at twelve, and about a quarter of that accumulation reverses by two years. An assessment of credit transmission that stops at one year records the effect at its maximum and overstates its durability, while one that looks only at two years misses the peak. Second, an easing surprise is not passed through one for one to new mortgage borrowers within the window studied. The mortgage margin over the deposit facility rate widens by about 0.05 percentage points one year after a typical surprise, and the two-year estimate is positive but imprecise. Third, the response is sectoral. Broad NFC and household credit stocks do not move with any precision after the same surprise that moves house-purchase lending growth, so a monitoring framework built on total credit would have recorded little, while the channel that responded most durably sat inside it. These implications are analytical rather than prescriptive.

The results also bear on the affordability debate, in a specific and limited way. \Cref{chap:introduction} set out the backdrop: real German house prices rose by about half between 2015 and 2021 while real compensation barely moved. Any account of that gap that runs through monetary policy has to pass through housing credit, and this thesis studies that step. Easing-signed surprises moved house-purchase lending growth at short and medium horizons and lending margins durably, while they reached the euro-area wage tracker weakly and briefly. That is the credit side of the mechanism the affordability debate presumes, established at monthly frequency with a high-frequency shock. It is not evidence on who obtained credit or on what happened to prices, and it says nothing about access being widened or narrowed for any group. Closing that loop would need three further pieces of evidence: house prices at a frequency that allows a monthly response, loan-level data on who borrowed and at what terms, and household balance-sheet data on who owned housing before the shocks arrived. Without them the thesis stops at the channel, and it says so.

For research, the results carry two lessons. A study that uses raw ECB target surprises as policy shocks should screen for information effects before it reads equity or expectation responses. In this sample 132 of the 208 announcements in the identification window carry a sign conflict or a possible information component, and the raw surprise produces a negative equity response that the screen helps to interpret. And work on pass-through at the monthly frequency gains from treating the margin over the policy rate as the measured object, not only the lending rate. The lending-rate levels in this thesis are estimated imprecisely, while the margin built from them is the most stable response in the whole set, with the same sign in all 42 sample-by-horizon cells, because netting out the policy floor removes the common movement that the rate levels share.

\section{Limitations}

The limits of the design set the limits of the claims. The shock is a high-frequency ECB policy-news surprise, not a purified structural monetary innovation. The best low-frequency bridge, the end-of-month change in the deposit facility rate, has real full-sample relevance but concentrated, event-driven strength, which \Cref{chap:robustness} documents in detail. That is why the shock is read as policy news rather than as an instrument for a single policy variable, and why information effects and event contamination stay inside the identification environment rather than being assumed away. It also separates this thesis from \textcite{GertlerKaradi2015} and \textcite{JarocinskiKaradi2020}, who use the surprises inside a VAR to recover structural shocks. The reduced-form route taken here is more transparent about what the surprise is, at the cost of not naming a treatment variable.

The aggregation boundary matters just as much. Quarterly house prices, compensation per employee, and the Bank Lending Survey are not interpolated into monthly responses. Household-level access, debt-service ratios, borrower heterogeneity, regional house prices, tenure status, and bank-level balance sheets are outside the empirical design, and so is any observed sequence from lending conditions to housing credit, as the previous sections made clear. The statistical caveats are the familiar ones. Many response variables, mixed signs, proxy dependence, and horizon dependence all call for caution in reading any single coefficient. The accumulated intervals depend, in addition, on how the dependence across horizons is handled. The HAC intervals treat the horizon estimates as independent, and the joint bootstraps that relax this widen the two-year housing-credit band enough to include zero. The thesis reports both and rests the persistence claim on the horizons where they agree.

Inside those boundaries the conclusion is still substantive. Under a reproducible monthly local-projection framework, ECB policy-news shocks show clearer evidence in house-purchase lending growth and lending spreads than in the euro-area wage tracker and broad productive-credit quantities, with the German industry wage bill as the exception on the compensation side. That is a disciplined result about heterogeneous macro-financial transmission. It is not a welfare or inequality claim.

\section{Future Research Agenda}
\label{sec:future-research}

Each limitation points to a specific extension. The aggregation boundary comes first. Bank-level and loan-level data would show whether the co-movement of margins and housing credit documented here reflects lender heterogeneity, funding-cost pass-through, margin adjustment, or credit-supply behavior, and would make an actual sequence from lending conditions to credit observable where monthly aggregates cannot. The household boundary comes second. Household-level credit, income, and balance-sheet data would allow incidence analysis across borrowers, renters, homeowners, income groups, and refinancing status, and would turn the comparative compensation result into a statement about who gained. The frequency boundary on house prices comes third. Regional housing data would show whether the housing-credit response is stronger where local housing constraints are already tight, and would connect the credit result to prices without interpolating quarterly series. The identification limits come last. A longer post-COVID sample, and a shock series purged of information effects along the lines of \textcite{JarocinskiKaradi2020}, would improve the regime analysis and allow the first-stage bridge to the policy rate to be tested again on a sample less dominated by a few events. Each direction relaxes one boundary of the present design without discarding its identification discipline.
